% ==============================================================================
% Section 5: Experiments
% ==============================================================================

\section{Experiments}
\label{sec:experiments}

\begin{table*}[!t]
\centering
\caption{
  \textbf{Confidence-gated ensemble, strict overlap@1 (\%).}
  \textbf{Native}: backbone baseline. \textbf{+ZwerGe}: margin-gated ensemble (\S\ref{sec:method:ensemble}), same cached run. $\Delta=$+ZwerGe$-$Native. \textbf{\%Fus.}: fraction routed to fusion. OSW-G = OSWorld-G (original); OSW-G-Ref = refusal-removed. Global $\tau{=}0.20$; 20/24 improve, 21/24 non-negative, macro-average $+0.46$\,pp.
}
\label{tab:main_ensemble}
\small
\setlength{\tabcolsep}{6pt}
\renewcommand{\arraystretch}{1.04}
\begin{tabular}{ll rrr r}
\toprule
\textbf{Backbone} & \textbf{Benchmark} & \textbf{Native} & \textbf{+ZwerGe} & \textbf{$\Delta$} & \textbf{\%Fus.} \\
 & & \textbf{ov@1} & \textbf{ov@1} & & \\
\midrule
\multirow{6}{*}{\shortstack[l]{GUI-Owl-7B~\citep{ye2025guiowl}\\{\scriptsize Qwen2.5-VL}}}
 & SS-Pro        & 55.22 & \textbf{56.10} & $+0.89$ & 7.5 \\
 & SS-v2         & 92.84 & \textbf{93.08} & $+0.24$ & 5.8 \\
 & OSW-G-Ref     & 79.61 & \textbf{80.39} & $+0.78$ & 25.7 \\
 & OSW-G         & 72.55 & \textbf{72.94} & $+0.39$ & 25.3 \\
 & MMBench-GUI   & 80.77 & \textbf{81.27} & $+0.50$ & 7.5 \\
 & UI-Vision     & 53.87 & \textbf{54.42} & $+0.55$ & 12.7 \\
\midrule
\multirow{6}{*}{\shortstack[l]{UI-TARS-1.5-7B~\citep{qin2025uitars}\\{\scriptsize Qwen2.5-VL}}}
 & SS-Pro        & 48.32 & \textbf{49.59} & $+1.27$ & 6.8 \\
 & SS-v2         & 91.35 & \textbf{91.66} & $+0.31$ & 5.7 \\
 & OSW-G-Ref     & 77.06 & 76.47 & $-0.59$ & 24.5 \\
 & OSW-G         & 64.51 & \textbf{65.10} & $+0.59$ & 20.8 \\
 & MMBench-GUI   & 77.46 & \textbf{77.88} & $+0.42$ & 7.6 \\
 & UI-Vision     & 27.12 & \textbf{28.45} & $+1.33$ & 7.6 \\
\midrule
\multirow{6}{*}{\shortstack[l]{GUI-Owl-1.5-8B~\citep{xu2026mobile}\\{\scriptsize Qwen3-VL}}}
 & SS-Pro        & 74.89 & \textbf{75.21} & $+0.32$ & 10.2 \\
 & SS-v2         & 94.10 & 94.02 & $-0.08$ & 9.9 \\
 & OSW-G-Ref     & 80.98 & \textbf{82.35} & $+1.37$ & 32.7 \\
 & OSW-G         & 78.43 & \textbf{78.63} & $+0.20$ & 30.4 \\
 & MMBench-GUI   & 83.58 & \textbf{84.17} & $+0.58$ & 9.9 \\
 & UI-Vision     & 43.94 & \textbf{44.97} & $+1.04$ & 14.2 \\
\midrule
\multirow{6}{*}{\shortstack[l]{UI-Venus-1.5-8B~\citep{team2026ui}\\{\scriptsize Qwen3-VL}}}
 & SS-Pro        & 70.52 & \textbf{70.65} & $+0.13$ & 7.7 \\
 & SS-v2         & 96.85 & 96.85 & $0.00$ & 8.6 \\
 & OSW-G-Ref     & 86.08 & \textbf{86.27} & $+0.20$ & 28.8 \\
 & OSW-G         & 81.76 & \textbf{82.35} & $+0.59$ & 25.5 \\
 & MMBench-GUI   & 89.04 & 88.95 & $-0.08$ & 8.8 \\
 & UI-Vision     & 52.18 & \textbf{52.33} & $+0.16$ & 11.8 \\
\bottomrule
\end{tabular}
\end{table*}

% ------------------------------------------------------------------------------
\subsection{Experimental Setup}
\label{sec:experiments:setup}
% ------------------------------------------------------------------------------

\paragraph{Benchmarks.}
We evaluate on six GUI grounding configurations: ScreenSpot-Pro~\citep{screenspotpro}, ScreenSpot-v2~\citep{Os-atlas}, the original and refusal-removed OSWorld-G-refine splits of OSWorld-G~\citep{xie2025scaling}, MMBench-GUI~\citep{mmbench}, and UI-Vision~\citep{ui-vision}, spanning desktop, web, and mobile.

\paragraph{Metrics.}
We report the strict single-point success rate (overlap@1); protocol and patch-mapping details are in the supplementary material.

\paragraph{Baselines.}
We retrofit four GUI agents: GUI-Owl-7B~\citep{ye2025guiowl} and UI-TARS-1.5-7B~\citep{qin2025uitars}, both built on the 28-layer Qwen2.5-VL, and GUI-Owl-1.5-8B~\citep{xu2026mobile} and UI-Venus-1.5-8B~\citep{team2026ui}, both built on the 36-layer Qwen3-VL. Each backbone is frozen throughout; only the probe and fusion heads are updated. The Native column is each backbone's own autoregressive accuracy on the same task set, so the comparison isolates the retrofit.

\paragraph{Implementation.}
Stage~1 trains 14 (Qwen2.5-VL) or 18 (Qwen3-VL) CrossAttn probes per model on 200k grounding samples for one epoch; Stage~2 trains the fusion head with a subset of active probes, 10 layers for the Qwen2.5-VL backbones and 12 for the Qwen3-VL backbones, for one epoch, with the backbone and the remaining probes frozen. The deployed decoder is the margin-gated ensemble (\S\ref{sec:method:ensemble}) with $\tau{=}0.20$.

% ------------------------------------------------------------------------------
\subsection{Main Results}
\label{sec:experiments:main}
% ------------------------------------------------------------------------------

Table~\ref{tab:main_ensemble} compares the confidence-gated ensemble against the native baseline under the identical strict success-rate protocol, with a single global threshold and no per-backbone tuning. Twenty of 24 cells improve and 21 of 24 are non-negative, for a macro-average gain of $+0.46$\,pp; the three negative cells are all within $0.6$\,pp of zero. The gate routes to the fusion point on 5.7--32.7\% of samples and otherwise defers to the native decoder; it fires most on OSWorld-G, where the native baseline has the most headroom, and least on ScreenSpot-v2, where native accuracy is near ceiling.

\textbf{Interpreting the gains.}
We read these small, mostly positive gains as evidence that intermediate-layer signals can complement native decoding, not as a new state-of-the-art result. The magnitude is small by construction: the gate defers to the native decoder on the majority of samples and overrides only when its posterior is sharply peaked. On the samples where the two disagree, the ensemble is net-positive in the large majority of cells, and a pooled paired test over all samples is decisive in its favor ($p{<}10^{-17}$); the gain is individually significant ($p{<}0.05$) on the six highest-headroom cells, and no cell is significantly worse than the baseline (per-cell table in the supplementary).

% ------------------------------------------------------------------------------
\subsection{Diagnosing the Gains: Target Geometry and Retrofit Headroom}
\label{sec:experiments:sspro}
% ------------------------------------------------------------------------------

We explain the two patterns in Table~\ref{tab:main_ensemble}, the largest headroom on ScreenSpot-Pro and the smallest gains on Qwen3-VL, with two diagnostics computed from the same cached runs.

\textbf{Why ScreenSpot-Pro is the hard case.}
Roughly three quarters of ScreenSpot-Pro targets are smaller than a single visual patch (see supplementary), so a patch-level posterior can at best identify the correct patch, not the sub-patch pixel. This is precisely why the retrofit is framed as a \emph{gate} rather than a replacement decoder: the intermediate posterior identifies \emph{which patch} contains the target, while sub-patch precision is left to the native serializer whenever the probe is uncertain.

\textbf{Retrofit headroom predicts the gate's behavior.}
Define the \emph{retrofit headroom} $h=\max_\ell\,\mathrm{ov@1}(p_\ell)-\mathrm{ov@1}(\text{native})$. When $h>0$, bypassing the serializer recovers real headroom; when $h\le0$, the native decoder has already closed that gap. A supplementary figure plots $h$ against the ungated fusion-only delta. Qwen2.5-VL backbones, whose native decoders leave more headroom, occupy the positive-headroom, positive-gain quadrant; Qwen3-VL backbones, whose native decoders approach the intermediate-layer ceiling, sit near or below zero. This is the regime the confidence gate targets: it defers to the native decoder whenever the fusion posterior of Eq.~\ref{eq:fusion} is diffuse, which is why most cells in Table~\ref{tab:main_ensemble} are non-negative.


